\section{Related Work}
\label{sec:related}

\paragraph{Length generalisation and positional encodings.} Transformers fit short
sequences and degrade sharply beyond them
\citep{anil2022exploring,deletang2023neural}. The explanations offered are mostly
computational: attention of fixed depth cannot represent certain length-dependent
functions \citep{hahn2020theoretical}; a task extrapolates when it admits a short
program over attention primitives \citep{weiss2021thinking,zhou2023what}; shallow
models learn shortcuts that fit the training lengths and break outside them
\citep{liu2023shortcuts}. Orthogonal to this is the choice of positional code---a
distinct vector per position \citep{vaswani2017attention}, or pairwise information
injected through learned relative representations \citep{shaw2018self}, bucketed
biases \citep{raffel2020exploring}, rotations \citep{su2024roformer} or a linear
distance penalty \citep{press2022train}---together with a family of schemes built for
the extrapolation regime: log-decaying biases \citep{chi2023dissecting},
resolution-maximising codes \citep{sun2023length}, randomised position sampling
\citep{ruoss2023randomized}, and priming with long examples
\citep{jelassi2023length}; \citet{zhao2024survey} survey the family. The headline
result is that no positional encoding beats five explicit alternatives across ten
tasks \citep{kazemnejad2023impact}, with follow-up work bounding how far that holds
and tracing the residual breakdown to unfocused attention \citep{wang2024length}.
Almost all of it is scored by exact match.

\paragraph{The EOS decision.} A separate literature studies the decision to predict an
end-of-sequence token at all. \citet{newman2020eos} show that models trained to predict
EOS extrapolate substantially worse than models that are not, trace the difference to
hidden states that stratify by position and then stall, and evaluate in the
oracle-length setting we adopt in \secref{sec:stop}. Their experiments vary the EOS
decision and hold the architecture fixed; ours varies the positional encoding and holds
the EOS decision fixed, then removes it---which is what connects the two literatures.
Small algorithmic models are the usual place to look for mechanism---grokking as a
transition from memorisation to a structured circuit
\citep{power2022grokking,nanda2023progress}, induction heads as a copying mechanism
\citep{olsson2022induction}---and we offer behaviour of that kind rather than a circuit
analysis: the profile in \figref{fig:profile} constrains what such an account must
explain, but does not localise the computation.
